\chapter{Implementation}

\section{Implementation Overview}

This chapter explains how \textbf{VisionSafe360} was implemented as a complete AI-based workplace safety monitoring system. Unlike the design chapters, this chapter focuses on the practical implementation: the models that were loaded, the rules used to convert detections into hazards, the backend endpoints that receive incidents, and the dashboard mechanisms used to display live safety information.

The implemented system monitors video streams and detects four main hazard categories: missing personal protective equipment (PPE), unsafe worker--forklift proximity, worker falls, and harmful ergonomic posture. These modules run in an edge AI pipeline and send structured events to a FastAPI backend. The backend stores incidents, publishes real-time updates through WebSockets, and exposes data to the React dashboard.

\section{Development Milestones}

The implementation was completed in five main milestones. The first milestone defined the hazard categories and translated them into measurable computer vision tasks. PPE violations were represented as missing-equipment detections, forklift risk was represented as worker--vehicle distance, falls were represented as temporal posture transitions, and ergonomics were represented as RULA/REBA scores.

The second milestone focused on dataset preparation and model selection. YOLO-based models were selected for object detection and pose estimation because they provide real-time inference and integrated tracking support through the Ultralytics framework \cite{ultralytics_yolov8}.

The third milestone implemented the edge AI modules. Each hazard module was developed as a separate analyzer so that the pipeline could enable or disable PPE, proximity, fall, and posture logic independently.

The fourth milestone implemented the backend and dashboard. FastAPI was used for authenticated API endpoints, PostgreSQL was used for persistent records, Redis supported real-time event delivery and worker orchestration, and React/TypeScript was used for the dashboard interface.

The final milestone integrated the system using REST, WebSocket streams, and Docker Compose so that the prototype could run as a connected multi-service application.

\section{System Architecture Implementation}

VisionSafe360 was implemented as a layered architecture consisting of video input, edge AI processing, backend services, and dashboard visualization.

The \textbf{video input layer} receives frames from CCTV cameras, webcams, IP cameras, RTSP streams, or uploaded video files. The stream handler keeps only the latest frame when necessary so that processing remains close to real time instead of building a large queue of old frames.

\begin{figure}[H]
\centering
\safeincludegraphics[width=0.95\linewidth,height=0.55\textheight,keepaspectratio]{images/video_input_pipeline.png}
\caption{Video input pipeline}
\end{figure}

The \textbf{edge AI processing layer} owns the GPU inference loop. It loads the pose, PPE, and proximity models, performs YOLO inference, applies ByteTrack tracking where needed, and passes detections to the hazard analyzers.

\begin{figure}[H]
\centering
\safeincludegraphics[width=0.95\linewidth,height=0.55\textheight,keepaspectratio]{images/edge_ai_pipeline.png}
\caption{Edge AI processing pipeline}
\end{figure}

The \textbf{backend layer} receives incidents from edge workers, validates requests, stores records, applies rate limits, and broadcasts incident updates to connected dashboard clients.

The \textbf{dashboard layer} provides live monitoring, incident management, alerts, reports, camera management, ergonomics analytics, and system configuration.

\begin{figure}[H]
\centering
\safeincludegraphics[width=0.95\linewidth,height=0.60\textheight,keepaspectratio]{images/fig_5_3_complete_architecture.png}
\caption{Complete VisionSafe360 architecture}
\end{figure}

\section{Technologies and Tools Used}

Table~\ref{tab:chapter5-technologies} summarizes the main tools used in the implementation. Version numbers were taken from the project requirements and package configuration files.

\begin{table}[H]
\centering
\footnotesize
\renewcommand{\arraystretch}{1.45}
\begin{tabular}{|p{0.2\linewidth}|p{0.16\linewidth}|p{0.28\linewidth}|p{0.25\linewidth}|}
\hline
\textbf{Tool} & \textbf{Version} & \textbf{Purpose} & \textbf{Justification} \\
\hline
Python & 3.11 & Edge AI and backend services & Strong AI and API ecosystem \\
\hline
Ultralytics YOLO & 8.4.16 & Detection, pose estimation, and tracking & Real-time model execution and ByteTrack integration \\
\hline
PyTorch & 2.10.0 + CUDA 12.8 & GPU inference & Efficient tensor execution on NVIDIA GPU \\
\hline
OpenCV & 4.13.0 / 4.10.0 & Frame handling and image processing & Mature video processing library \\
\hline
FastAPI & 0.115.0 & Backend REST API and WebSocket routing & Fast asynchronous API framework \cite{fastapi_docs} \\
\hline
PostgreSQL & 16 & Persistent database & Reliable relational storage \cite{postgres_docs} \\
\hline
SQLAlchemy & 2.0.35 & ORM/database layer & Structured database models and migrations \\
\hline
Redis & 5.0.8 client, Redis 7 services & Queues, rate counters, event streams, latest-frame cache & Low-latency shared state and Pub/Sub \cite{redis_docs} \\
\hline
React & 19.2.1 & Dashboard UI & Component-based interface development \cite{react_docs} \\
\hline
TypeScript & 5.8.2 & Dashboard application logic & Type safety for frontend data models \\
\hline
Tailwind CSS & 3.4.4 & Dashboard styling & Consistent utility-first UI styling \\
\hline
Docker Compose & Compose stack & Multi-service deployment & Reproducible backend, dashboard, database, Redis, and worker services \cite{docker_compose_docs} \\
\hline
\end{tabular}
\caption{\textbf{Technologies and tools used in VisionSafe360}}
\label{tab:chapter5-technologies}
\end{table}

\section{Hardware Environment}

The project was developed and tested on a laptop with GPU acceleration. The hardware was sufficient for running YOLO inference, backend services, and the dashboard during prototype testing.

\begin{table}[H]
\centering
\renewcommand{\arraystretch}{1.55}
\begin{tabular}{|>{\centering\arraybackslash}p{0.3\linewidth}|>{\centering\arraybackslash}p{0.58\linewidth}|}
\hline
\textbf{Component} & \textbf{Specification} \\
\hline
\textbf{CPU} & Intel Core i5-13420H \\
\hline
\textbf{GPU} & NVIDIA GeForce RTX 4050 Laptop GPU \\
\hline
\textbf{RAM} & 16 GB \\
\hline
\textbf{Operating System} & Windows 11 development environment \\
\hline
\end{tabular}
\caption{\textbf{Hardware environment used for development and testing}}
\label{tab:chapter5-hardware}
\end{table}

\section{Module Implementation}

The edge AI implementation is divided into independent modules. Each module receives detections or pose keypoints and returns a list of structured hazard events. This modular design makes the system easier to test and allows specific hazards to be enabled through runtime profiles.

\subsection{PPE Detection}

The PPE detection module monitors whether workers are wearing required equipment such as helmets, safety vests, gloves, and safety shoes. The module uses a YOLO model trained on SH17, a public PPE dataset containing 8,099 images and 75,994 labeled instances across 17 classes \cite{ahmad2024sh17}. In the implementation, SH17-style labels are normalized into application labels such as \texttt{helmet\_on}, \texttt{helmet\_off}, \texttt{safety\_vest}, \texttt{vest\_off}, \texttt{gloves\_off}, and \texttt{bare\_foot}.

\begin{table}[H]
\centering
\renewcommand{\arraystretch}{1.45}
\begin{tabular}{|>{\centering\arraybackslash}p{0.32\linewidth}|>{\centering\arraybackslash}p{0.56\linewidth}|}
\hline
\textbf{Item} & \textbf{Value} \\
\hline
\textbf{Dataset} & SH17 PPE Dataset \\
\hline
\textbf{Images / Instances} & 8,099 images / 75,994 instances \\
\hline
\textbf{Detected PPE Classes} & Helmet, safety vest, gloves, shoes, mask, glasses, face guard, ear protection, suit, tool, and related missing-equipment classes \\
\hline
\textbf{Input Resolution} & 640$\times$640 \\
\hline
\textbf{Model Family} & YOLO detector \\
\hline
\textbf{Persistence Rule} & Missing PPE must persist before an incident is emitted \\
\hline
\end{tabular}
\caption{\textbf{PPE dataset and implementation configuration}}
\label{tab:chapter5-ppe-dataset}
\end{table}

Listing~\ref{lst:ppe-inference} shows the core PPE inference pattern. The model is loaded once, then each frame is passed to YOLO with confidence, IoU, and maximum-detection limits.

\begin{lstlisting}[language=Python,caption={PPE model inference on one frame},label={lst:ppe-inference}]
from ultralytics import YOLO

ppe_model = YOLO("weights/ppe/best_ppe.pt")

results = ppe_model.predict(
    frame,
    device="cuda:0",
    imgsz=640,
    conf=0.30,
    iou=0.50,
    max_det=100,
    verbose=False,
)[0]

detections = []
for box in results.boxes:
    cls_id = int(box.cls[0])
    name = ppe_model.names[cls_id].lower()
    detections.append((name, float(box.conf[0]), box.xyxy[0]))
\end{lstlisting}

The analyzer converts missing-equipment detections into hazard events. For example, \texttt{helmet\_off} is converted into a high-severity missing-helmet event, while \texttt{gloves\_off} is treated as medium severity. The event aggregator then applies persistence and cooldown logic to reduce false positives caused by short occlusions.

\begin{figure}[H]
\centering
\safeincludegraphics[width=0.95\linewidth,height=0.58\textheight,keepaspectratio]{images/Vest, Helmet, Gloves detection.png}
\caption{PPE detection output}
\end{figure}

\subsection{Forklift Proximity Detection}

The forklift proximity module detects unsafe distances between workers and forklifts. The implementation uses a YOLO detector for forklifts and persons, then calculates Euclidean distance between bounding box centroids in pixel space. The current prototype uses pixel thresholds because monocular CCTV footage does not provide direct depth information. Camera calibration can later replace these pixel thresholds with meter-based distances.

\begin{table}[H]
\centering
\renewcommand{\arraystretch}{1.45}
\begin{tabular}{|>{\centering\arraybackslash}p{0.34\linewidth}|>{\centering\arraybackslash}p{0.54\linewidth}|}
\hline
\textbf{Item} & \textbf{Value} \\
\hline
\textbf{Model} & YOLOv8m forklift/person detector \\
\hline
\textbf{Training Images} & 2,468 \\
\hline
\textbf{Validation Images} & 500 \\
\hline
\textbf{Epochs / Batch Size} & 200 epochs / batch size 16 \\
\hline
\textbf{Danger Threshold} & 80 pixels \\
\hline
\textbf{Warning Threshold} & 140 pixels \\
\hline
\textbf{Alert Persistence} & Unsafe proximity must persist for approximately one second \\
\hline
\end{tabular}
\caption{\textbf{Forklift proximity model and rule configuration}}
\label{tab:chapter5-forklift-model}
\end{table}

Listing~\ref{lst:proximity} shows the core rule. The nearest forklift is selected for each person, and severity is assigned based on the configured danger and warning radii.

\begin{lstlisting}[language=Python,caption={Worker--forklift distance check},label={lst:proximity}]
def center(bbox):
    x1, y1, x2, y2 = bbox
    return ((x1 + x2) / 2.0, (y1 + y2) / 2.0)

def distance_px(a, b):
    ax, ay = center(a)
    bx, by = center(b)
    return ((ax - bx) ** 2 + (ay - by) ** 2) ** 0.5

for person in persons:
    dist_px, forklift = min(
        (distance_px(person.bbox, f.bbox), f) for f in forklifts
    )
    if dist_px <= 80:
        emit_event("forklift_proximity_danger", "High", person)
    elif dist_px <= 140:
        emit_event("forklift_proximity_warning", "Medium", person)
\end{lstlisting}

\begin{figure}[H]
\centering
\safeincludegraphics[width=0.95\linewidth,height=0.58\textheight,keepaspectratio]{images/Forklift and worker detection.png}
\caption{Forklift detection output}
\end{figure}

\subsection{Fall Detection}

The fall detection module uses a pose-aware rule engine rather than a heavy sequence model. This decision reduced computational cost and made the system easier to explain and tune. Each tracked person has an independent state machine with three implementation states: \texttt{normal}, \texttt{candidate}, and \texttt{confirmed}. A recovered posture returns the track to \texttt{normal}.

The system evaluates bounding box aspect ratio, downward centroid velocity, hip position from COCO keypoints, and immobility. A fall is confirmed only after the candidate posture persists for 2.5 seconds, reducing false alarms from bending, crouching, or short occlusions.

\begin{lstlisting}[language=Python,caption={Fall state transition logic},label={lst:fall-state}]
if state == "normal" and triggered:
    state = "candidate"
    candidate_since = timestamp

elif state == "candidate":
    recovered = hip_ratio > 0.6 if hip_ratio else aspect_ratio < 0.8
    if recovered:
        state = "normal"
    elif timestamp - candidate_since >= 2.5:
        state = "confirmed"
        emit_event(
            "fall_confirmed",
            severity="Critical",
            metadata={"aspect_ratio": aspect_ratio,
                      "hip_ratio": hip_ratio},
        )

elif state == "confirmed" and recovered:
    state = "normal"
\end{lstlisting}

\begin{figure}[H]
\centering
\safeincludegraphics[width=0.95\linewidth,height=0.58\textheight,keepaspectratio]{images/fig_5_7_fall_state_machine.png}
\caption{Fall detection state machine}
\end{figure}

\begin{figure}[H]
\centering
\safeincludegraphics[width=0.95\linewidth,height=0.58\textheight,keepaspectratio]{images/Fallen worker .png}
\caption{Fall detection alert}
\end{figure}

\subsection{Ergonomic Posture Analysis}

The ergonomic posture module computes body angles from COCO-17 pose keypoints and feeds them into RULA and REBA scoring functions \cite{mcatamney1993rula,hignett2000reba}. Joint angles are derived from keypoint groups. For example, trunk flexion is computed from the midpoint between both hips and the midpoint between both shoulders, while upper-arm and lower-arm angles use shoulder, elbow, wrist, and hip keypoints.

\begin{table}[H]
\centering
\renewcommand{\arraystretch}{1.55}
\begin{tabular}{|>{\centering\arraybackslash}p{0.34\linewidth}|>{\centering\arraybackslash}p{0.54\linewidth}|}
\hline
\textbf{Metric} & \textbf{Threshold} \\
\hline
\textbf{RULA Alert} & $\geq$ 5 \\
\hline
\textbf{REBA Alert} & $\geq$ 7 \\
\hline
\textbf{Critical Risk} & RULA = 7 or REBA $\geq$ 11 \\
\hline
\textbf{Sustained Poor Posture} & 4 seconds before event emission \\
\hline
\textbf{Smoothing} & Exponential moving average with alpha = 0.6 \\
\hline
\end{tabular}
\caption{\textbf{Ergonomic posture alert thresholds}}
\label{tab:chapter5-ergonomic-thresholds}
\end{table}

Listing~\ref{lst:angle} shows the angle calculation used by the ergonomics implementation. The function computes the angle at a joint by comparing the vectors connected to that joint.

\begin{lstlisting}[language=Python,caption={Joint angle calculation for RULA/REBA},label={lst:angle}]
def angle_between(a, b, c):
    ba = a - b
    bc = c - b
    norm_ba = np.linalg.norm(ba)
    norm_bc = np.linalg.norm(bc)
    if norm_ba < 1e-6 or norm_bc < 1e-6:
        return None
    cos_angle = np.dot(ba, bc) / (norm_ba * norm_bc)
    cos_angle = np.clip(cos_angle, -1.0, 1.0)
    return float(np.degrees(np.arccos(cos_angle)))

angles = compute_angles(smoothed_keypoints, confidences)
rula_score = compute_rula(angles)["final_score"]
reba_score = compute_reba(angles)["final_score"]
\end{lstlisting}

\begin{figure}[H]
\centering
\safeincludegraphics[width=0.95\linewidth,height=0.58\textheight,keepaspectratio]{images/posture skeleton.png}
\caption{Posture skeleton detection}
\end{figure}

\subsection{Tracking and Event Persistence}

Tracking is used to maintain worker identity across frames and prevent duplicate incidents. The implementation uses Ultralytics \texttt{model.track(...)} with ByteTrack. The tracker provides persistent \texttt{track\_id} values that are passed to the PPE, proximity, fall, and ergonomics modules.

\begin{lstlisting}[language=Python,caption={ByteTrack initialization and ID assignment},label={lst:tracking}]
results = pose_model.track(
    frame,
    device="cuda:0",
    imgsz=640,
    conf=0.38,
    iou=0.50,
    tracker="bytetrack.yaml",
    persist=True,
    verbose=False,
)[0]

for box in results.boxes:
    track_id = int(box.id[0]) if box.id is not None else None
    detections.append(Detection(
        class_name="person",
        confidence=float(box.conf[0]),
        bbox=tuple(map(int, box.xyxy[0])),
        track_id=track_id,
    ))
\end{lstlisting}

Event persistence is handled by temporal rules. PPE events must remain visible before being reported, proximity events are emitted only when the threshold condition persists, fall events require a candidate confirmation window, and posture events require sustained high ergonomic risk. This design prevents single-frame errors from becoming incidents.

\section{Backend Implementation}

The backend was developed using FastAPI \cite{fastapi_docs}. It acts as the central communication layer for incidents, alerts, analytics, camera configuration, authentication, jobs, and real-time streams.

\begin{table}[H]
\centering
\footnotesize
\renewcommand{\arraystretch}{1.35}
\begin{tabular}{|p{0.24\linewidth}|p{0.24\linewidth}|p{0.41\linewidth}|}
\hline
\textbf{Endpoint Group} & \textbf{Method} & \textbf{Purpose} \\
\hline
\texttt{/api/incidents} & GET/POST/PATCH/DELETE & List, create, update, and delete incident records \\
\hline
\texttt{/api/alerts} & GET/POST/PATCH & Manage alert records and acknowledgement workflow \\
\hline
\texttt{/api/analytics/*} & GET & Provide time-series, severity, zone, and dashboard analytics \\
\hline
\texttt{/api/cameras} & GET/POST/PATCH/DELETE & Manage camera sources and metadata \\
\hline
\texttt{/api/ergonomics} & GET/POST & Store and query ergonomic score samples \\
\hline
\texttt{/api/jobs} & GET/POST & Start and monitor AI processing jobs \\
\hline
\texttt{/ws/incidents} & WebSocket & Push incident updates to dashboard clients \\
\hline
\texttt{/ws/stream/\{camera\_id\}} & WebSocket & Stream annotated AI frames to the dashboard \\
\hline
\end{tabular}
\caption{\textbf{Backend endpoint overview}}
\label{tab:chapter5-api}
\end{table}

The database layer uses SQLAlchemy models for incidents, alerts, cameras, users, notifications, ergonomic records, and system configuration. Alembic migrations are included to create and evolve the PostgreSQL schema. Redis is used for rate limiting, event Pub/Sub, worker queues, orchestration state, and latest-frame streaming.

Listing~\ref{lst:incident-endpoint} shows the core incident ingestion endpoint. After creating a database record, the backend broadcasts an \texttt{incident\_created} event to connected dashboard clients.

\begin{lstlisting}[language=Python,caption={FastAPI incident ingestion endpoint},label={lst:incident-endpoint}]
@router.post("", response_model=IncidentOut, status_code=201)
async def create_incident(payload: IncidentCreate,
                          request: Request,
                          db: Session = Depends(get_db),
                          x_source_id: str | None = Header(default=None)):
    source_id = (x_source_id or payload.zone or "unknown").strip()
    allowed, retry_after = rate_limit_service.check_and_consume(source_id)
    if not allowed:
        raise HTTPException(status_code=429,
                            detail="Incident ingestion rate limit exceeded")

    incident = IncidentService.create(db, payload)
    await incident_ws_manager.broadcast({
        "type": "incident_created",
        "timestamp": datetime.now(timezone.utc).isoformat(),
        "incident": serialize_incident(incident),
    })
    return incident
\end{lstlisting}

\begin{figure}[H]
\centering
\safeincludegraphics[width=0.95\linewidth,height=0.58\textheight,keepaspectratio]{images/Fast Api swagger.png}
\caption{Backend API documentation}
\end{figure}

\section{Dashboard Implementation}

The dashboard was implemented using React and TypeScript \cite{react_docs}. It gives supervisors and safety teams access to live video, active incidents, historical reports, camera management, and ergonomic analytics.

The main dashboard pages are Dashboard Home, Live Monitoring, Alerts, Incident Management, Reports, Ergonomics Analytics, Camera Management, System Configuration, User Management, and System Health. The Live Monitoring page is the most integration-heavy view because it connects to \texttt{/ws/incidents} for real-time incident updates and to stream WebSocket endpoints for annotated AI frames.

\begin{lstlisting}[language=JavaScript,caption={Dashboard WebSocket incident feed},label={lst:dashboard-ws}]
const ws = new WebSocket(`${WS_BASE_URL}/ws/incidents?token=${token}`);

ws.onmessage = (event) => {
  const payload = JSON.parse(event.data);
  if (payload?.type === "incident_created" && payload?.incident) {
    upsertIncident(normalizeIncident(payload.incident));
  }
};

ws.onclose = () => {
  setTimeout(connectIncidentFeed, 3000);
};
\end{lstlisting}

\begin{figure}[H]
\centering
\safeincludegraphics[width=0.95\linewidth,height=0.58\textheight,keepaspectratio]{images/img_019.png}
\caption{Main dashboard}
\end{figure}

\begin{figure}[H]
\centering
\safeincludegraphics[width=0.95\linewidth,height=0.58\textheight,keepaspectratio]{images/Live Monitoring .png}
\caption{Live monitoring page}
\end{figure}

\begin{figure}[H]
\centering
\safeincludegraphics[width=0.95\linewidth,height=0.58\textheight,keepaspectratio]{images/img_021.png}
\caption{Incident management page}
\end{figure}

\begin{figure}[H]
\centering
\safeincludegraphics[width=0.95\linewidth,height=0.58\textheight,keepaspectratio]{images/img_022.png}
\caption{Ergonomics dashboard}
\end{figure}

\section{System Integration}

The complete VisionSafe360 prototype integrates the edge AI worker, backend, database, Redis services, dashboard, and media relay.

The edge AI worker converts each hazard into a \texttt{HazardEvent}. The backend client serializes this event into the backend incident schema and sends it through a REST \texttt{POST} request to \texttt{/api/incidents}. If the backend is unavailable, the payload is stored in a local SQLite offline queue and retried later.

After the backend stores the incident in PostgreSQL, it publishes a WebSocket event to all connected dashboard clients. The dashboard receives the event immediately and inserts or updates the incident in the live feed.

Docker Compose ties the services together. The stack includes PostgreSQL 16, multiple Redis 7 services, the FastAPI backend, an RQ worker for AI orchestration, the React dashboard served through Nginx, MediaMTX for RTSP/HLS/WebRTC relay, and pgAdmin for database inspection.

\begin{lstlisting}[language=Python,caption={Edge AI event delivery to backend},label={lst:backend-client}]
payload = event_to_payload(hazard_event)
headers = {"Content-Type": "application/json"}
if auth_token:
    headers["Authorization"] = f"Bearer {auth_token}"

response = session.post(
    "http://backend:8000/api/incidents",
    json=payload,
    timeout=5.0,
    headers=headers,
)

if response.status_code >= 300:
    enqueue_offline(payload)
\end{lstlisting}

\begin{lstlisting}[language={},caption={Docker Compose service integration excerpt},label={lst:compose}]
backend:
  ports: ["8000:8000"]
  depends_on: [db, redis-queue, redis-state, redis-events]

worker:
  command: python worker.py
  environment:
    VISIONSAFE_BACKEND_URL: http://backend:8000

dashboard:
  ports: ["80:80"]
  depends_on: [backend]
\end{lstlisting}

The resulting data flow is:

\begin{enumerate}
\item Camera or video file provides frames to the edge AI worker.
\item YOLO models detect people, PPE, forklifts, and pose keypoints.
\item Hazard analyzers create structured events with severity and metadata.
\item The edge client posts incidents to the backend REST API.
\item The backend stores records in PostgreSQL and publishes WebSocket updates.
\item The dashboard displays live incidents, annotated streams, reports, and ergonomic trends.
\end{enumerate}

\section{Challenges Encountered}

Table~\ref{tab:chapter5-challenges} summarizes the main implementation challenges and the solutions applied.

\begin{table}[H]
\centering
\footnotesize
\renewcommand{\arraystretch}{1.45}
\begin{tabular}{|p{0.25\linewidth}|p{0.3\linewidth}|p{0.34\linewidth}|}
\hline
\textbf{Challenge} & \textbf{Root Cause} & \textbf{Solution Implemented} \\
\hline
False-positive PPE alerts & Temporary occlusion, camera angle, or partial worker visibility & Persistence and cooldown logic before emitting incidents \\
\hline
Lighting variation & Low illumination reduced detector confidence & Confidence threshold tuning and model-specific preprocessing \\
\hline
Pose jitter & Keypoints moved slightly between frames even for stable posture & Exponential moving average smoothing before RULA/REBA scoring \\
\hline
Depth ambiguity in proximity detection & Single RGB cameras do not directly provide real-world distance & Pixel-distance thresholds with support for future camera calibration \\
\hline
Duplicate incident creation & Same track may remain hazardous for several frames & Track IDs, aggregation windows, cooldowns, and backend duplicate checks \\
\hline
Backend outages during live inference & Network or service failures can occur while processing video & Offline SQLite queue in the edge backend client \\
\hline
\end{tabular}
\caption{\textbf{Implementation challenges and solutions}}
\label{tab:chapter5-challenges}
\end{table}

\section{Summary}

This chapter presented the implementation of VisionSafe360 as a working AI safety monitoring prototype. The PPE module uses YOLO-based detection and persistence rules for missing equipment. The proximity module measures worker--forklift centroid distance with 80-pixel danger and 140-pixel warning thresholds. The fall module confirms incidents after a 2.5-second candidate window using aspect ratio, hip position, downward motion, and immobility checks. The ergonomics module computes RULA and REBA scores from COCO-17 keypoints and emits posture alerts after sustained risk.

The backend and dashboard complete the system by storing incidents, broadcasting WebSocket updates, displaying live monitoring views, and supporting Docker-based deployment. Together, these components show how VisionSafe360 moves from a design concept into an integrated workplace safety monitoring system.
